\documentclass[12pt]{article}
\usepackage[margin=1in]{geometry}
\usepackage{ulem}
\title{Text Formatting Showcase}
\author{Format Tester}
\begin{document}
\maketitle

\section{Basic Formatting}
This paragraph contains \textbf{bold text}, \textit{italic text}, and
\underline{underlined text}. We can also combine them: \textbf{\textit{bold
italic text}} is used for strong emphasis.

\section{Monospace and Code}
Inline code like \texttt{print("hello world")} should be wrapped in backticks.
Variable names like \texttt{max\_value} and \texttt{user\_input} should preserve
underscores correctly.

\section{Font Sizes}
{\tiny This is tiny text.}
{\small This is small text.}
{\normalsize This is normal text.}
{\large This is large text.}
{\Large This is Large text.}
{\LARGE This is LARGE text.}

\section{Text Styles}
Here is \textsc{Small Caps Text} and here is \textsf{sans-serif text}.
We also have \textsl{slanted text} which is similar to italic.

\section{Mixed Formatting Paragraph}
In this paragraph, we mix \textbf{bold words} with \textit{italic phrases} and
some \texttt{code snippets} all in the same sentence. This tests whether the
converter can correctly identify formatting spans within a single line of text
and apply the appropriate Markdown markers without overlapping or nesting errors.
\end{document}